%! AP Statistics — Unit 1, Lesson 1.8 — Homework
\documentclass[10pt]{article}
\usepackage{apstats-article}
\usepackage{apstats-boxes}

%=============================================================================
\begin{document}

\pageheader{Unit 1, Lesson 1.8}{Homework: Boxplots}
\namedateperiod

\vspace{0.10in}

\begin{tcolorbox}[colback=sky, colframe=navy, boxrule=0.8pt, arc=2mm,
    left=3mm, right=3mm, top=1.5mm, bottom=1.5mm]
\small Show all calculations. Interpret statistics in context. Use complete
sentences for AP-style questions.
\end{tcolorbox}

\vspace{0.12in}

% ════════════════════════════════════════════════════════════════════════════
% PART A — Five-number summary and outlier check
% ════════════════════════════════════════════════════════════════════════════
\begin{blurbbox}[Part A \quad Five-Number Summary \& Outlier Detection]{navylight}
\small

A physical education teacher recorded the number of push-ups completed by
12 students in one minute. The sorted data are:
\[8,\quad 12,\quad 15,\quad 18,\quad 20,\quad 22,\quad 25,\quad 28,\quad 30,\quad 33,\quad 36,\quad 61\]

\vspace{0.10in}

\begin{enumerate}[leftmargin=1.5em, itemsep=8pt]

  \item Find the five-number summary. Show your work for $Q_1$ and $Q_3$.\\[5pt]
        Lower half (6 values): \blank{8cm}\\[4pt]
        $Q_1 = \blank{2cm}$ \quad Median $= \blank{2cm}$ \quad $Q_3 = \blank{2cm}$\\[4pt]
        $\min = \blank{1.5cm}$ \quad $\max = \blank{1.5cm}$ \quad
        IQR $= Q_3 - Q_1 = \blank{2cm}$

  \item Apply the $1.5 \times \text{IQR}$ fence rule. Is 61 an outlier?\\[5pt]
        Lower fence: \blank{6cm}\\[4pt]
        Upper fence: \blank{6cm}\\[4pt]
        Is 61 an outlier? \blank{2cm} \quad Justify: \blank{6cm}

  \item If this were displayed as a \textbf{modified boxplot}, where would the
        right whisker end? \blank{3cm} push-ups. \\[4pt]
        What would the outlier dot represent? \blank{8cm}

\end{enumerate}

\end{blurbbox}

\vspace{0.12in}

% ════════════════════════════════════════════════════════════════════════════
% PART B — Read and interpret a boxplot
% ════════════════════════════════════════════════════════════════════════════
\begin{blurbbox}[Part B \quad Read and Interpret a Boxplot]{greenacc}
\small

The modified boxplot below displays the heights (cm) of seedlings in a greenhouse
after two weeks of growth.

\vspace{0.08in}
\begin{center}
\begin{tikzpicture}[scale=1.0]
  % Axis: 0–50, scale v*0.16; ticks every 5
  \draw[thick] (0.3,0) -- (8.5,0);
  \foreach \v/\lbl in {0/0, 0.80/5, 1.60/10, 2.40/15, 3.20/20,
                        4.00/25, 4.80/30, 5.60/35, 6.40/40, 7.20/45, 8.00/50} {
    \draw (\v,-0.12) -- (\v,0.12);
    \node[below, font=\small] at (\v,-0.12) {\lbl};
  }
  \node[below, font=\small\itshape] at (4.0,-0.56) {Seedling height (cm)};

  % scale: v*0.16
  % min=8 (x=1.28), Q1=12 (x=1.92), M=18 (x=2.88),
  % Q3=24 (x=3.84), last non-outlier=36 (x=5.76), outlier=46 (x=7.36)
  % IQR=12, upper fence=24+18=42; 46>42 so outlier confirmed
  \draw[thick, navy] (1.28, 0.38) -- (1.92, 0.38);
  \draw[thick, navy] (1.92, 0.10) rectangle (3.84, 0.66);
  \draw[thick, navy] (2.88, 0.10) -- (2.88, 0.66);
  \draw[thick, navy] (3.84, 0.38) -- (5.76, 0.38);
  \draw[thick, navy] (1.28, 0.22) -- (1.28, 0.54);
  \draw[thick, navy] (5.76, 0.22) -- (5.76, 0.54);
  \filldraw[navy] (7.36, 0.38) circle (0.11);
\end{tikzpicture}
\end{center}

\vspace{0.08in}

\begin{enumerate}[leftmargin=1.5em, itemsep=8pt]

  \item Read the five-number summary from the boxplot.\\[5pt]
        $\min = \blank{2cm}$ cm \quad $Q_1 = \blank{2cm}$ cm \quad
        $M = \blank{2cm}$ cm \quad $Q_3 = \blank{2cm}$ cm \quad
        right whisker end $= \blank{2cm}$ cm\\[4pt]
        The outlier dot is at approximately \blank{2cm} cm.\\[4pt]
        IQR $= Q_3 - Q_1 = \blank{2cm}$ cm

  \item Verify that the outlier dot (46 cm) exceeds the upper fence.
        Show your work.\\[5pt]
        Upper fence $= Q_3 + 1.5(\text{IQR}) = \blank{7cm}$\\[4pt]
        Is 46 cm an outlier? \blank{2cm} \quad Justify: \blank{5cm}

  \item Is the distribution skewed left, symmetric, or skewed right?
        Justify your answer using the boxplot.\\[5pt]
        \writeline\\[1pt]\writeline

\end{enumerate}

\end{blurbbox}

\vspace{0.12in}

% ════════════════════════════════════════════════════════════════════════════
% PART C — AP-Style
% ════════════════════════════════════════════════════════════════════════════
\begin{blurbbox}[Part C \quad AP-Style Free Response]{redacc}
\small

Using the seedling boxplot from Part B, write a complete description of the
distribution. Address \textbf{all four SOCS components} in context.
Use correct statistical vocabulary.

\vspace{0.08in}
\writeline\\[1pt]\writeline\\[1pt]\writeline\\[1pt]
\writeline\\[1pt]\writeline\\[1pt]\writeline

\end{blurbbox}

\end{document}
